%%%%%%%%%%  scan idx 178  |  folio 165  %%%%%%%%%%
d'où la
première assertion. La seconde en est un cas particulier, puisque, pour $G\in\Ob\mathcal C_{0S}$, i.e.\ $G$ libre et de type fini, et $E\in\Ob\mathcal C'_S$, on a $\underline{\Hom}(G,E)\in\Ob\mathcal C'_S$.

b) La première assertion est triviale. Prouvons maintenant, dans le cas (i) ou (ii), la pleine fidélité de $\Gamma_S:\mathcal C'_S\longrightarrow\Mod(A(S))$. Soient $F,G\in\Ob\mathcal C'_S$, et soit
\[
A_S^{(I)}\longrightarrow A_S^{(J)}\longrightarrow F\longrightarrow0
\tag{$*$}
\]
une présentation globale de $F$ (avec $I$ et $J$ finis dans le cas (i)). Appliquant à $(*)$ le foncteur $\Gamma_S$, qui est exact sur $\mathcal C'_S$ et commute, dans le cas (ii), aux limites inductives filtrantes, on obtient une suite exacte
\[
A(S)^{(I)}\longrightarrow A(S)^{(J)}\longrightarrow\Gamma_S(F)\longrightarrow0.
\tag{$**$}
\]
D'autre part, appliquant à $(*)$ (resp. $(**)$) le foncteur exact à gauche $\operatorname{Hom}_{A_S}(.,G)$ (resp. $\operatorname{Hom}_{A(S)}(\cdot,\Gamma_S(G))$), on obtient un diagramme commutatif où les lignes sont exactes :
\[
\begin{tikzcd}[column sep=small]
0 \arrow[r] & \operatorname{Hom}_{A_S}(F,G) \arrow[r] \arrow[d] &
\operatorname{Hom}_{A_S}(A_S^{(J)},G) \arrow[r] \arrow[d] &
\operatorname{Hom}_{A_S}(A_S^{(I)},G) \arrow[d] \\
0 \arrow[r] & \operatorname{Hom}_{A(S)}(F,G) \arrow[r] &
\operatorname{Hom}_{A(S)}(A(S)^{(J)},\Gamma_S(G)) \arrow[r] &
\operatorname{Hom}_{A(S)}(A(S)^{(I)},\Gamma_S(G)).
\end{tikzcd}
\]
Les deux flèches verticales de droite sont visiblement des isomorphismes il en est donc de même de la flèche verticale de gauche, ce qui démontre la pleine fidélité de $\Gamma_S:\mathcal C'_S\longrightarrow\Mod(A(S))$. La caractérisation de l'image essentielle, dans le cas (i) ou (ii), est immédiate et laissée au lecteur.

\textbf{Lemme 2.1.3.} Soient $A$ et $B$ des catégories abéliennes, et $\phi:A\to B$ un foncteur exact.
\begin{enumerate}[label=\alph*)]
\item Pour $M\in\Ob A$, les conditions suivantes sont équivalentes :
\begin{enumerate}[label=(\roman*)]
\item pour tout $L\in\Ob A$, la flèche canonique (induite par $\phi$)
\[
\Ext^n_{A}(L,M)\longrightarrow\Ext^n_{B}(\phi(L),\phi(M))
\]
est un isomorphisme pour tout $n$ ;
\item pour tout $L\in\Ob A$, la flèche canonique
